% SHARD-ID: AQM-95-HAAH-SHEAF-SITE-SPECTRUM-ORIENTATION
% SHARD-TITLE: Haah Sheaf, Site, and Spectrum Orientation
% SHARD-SUMMARY: Records a speculative orientation dictionary separating Haah physical sites from the spectrum of the translation algebra.
% SHARD-SUMMARY: Recasts Haah modules as sheaves over \(\Spec\mathbb F_p[\Lambda]\) while keeping qubits over the translation site object.
% SHARD-SUMMARY: Illustrates the distinction with infinite lattices, finite periodic quotients, split group algebras, and modular thickenings.
% SHARD-KEYWORDS: Haah, sheaf, translation algebra, site object, Spec, group algebra, module, finite periodic quotient

\section{Haah Sheaf, Site, and Spectrum Orientation}
\label{sec:haah-sheaf-site-spectrum-orientation}

\subsection{Status}

\textbf{Status: speculative orientation synthesis.}  This shard records a
working dictionary prompted by the question: if Haah's Laurent-polynomial
formalism is secretly sheaf-theoretic, what geometric object is underneath
the qubits?  It adds no new source, convention, run artifact, or theorem.  The
source-local inputs are AQM-91, AQM-92, AQM-94, and conventions
\texttt{(az)}--\texttt{(ba)}.  In particular, AQM-94 separates pre-Haah data,
scheme Haah data, finite Weyl-Haah data, and translation Haah data
(AQM-94, lines 45--70), and it warns that the group basis of
\(\mathbb F_p[\Lambda]\) is not the Zariski point set of
\(\Spec\mathbb F_p[\Lambda]\) (AQM-94, lines 114--136).

\subsection{The First Answer}

The cleanest first answer is:
\[
  \text{the qubits live over a translation site object, not over }
  \Spec\mathbb F_p[\Lambda].
\]
For an infinite translation-invariant model, the site object is
\[
  X_{\rm site}=\Lambda\times\{1,\ldots,q\},
\]
where \(\Lambda\) is the translation group and \(q\) is the number of
qubits or qudits per unit cell.  For a finite periodic model, replace
\(\Lambda\) by a finite quotient \(G\):
\[
  X_{\rm site}=G\times\{1,\ldots,q\}.
\]
A finite-support Pauli label is then a section with finite support over this
discrete site object.  Algebraically, for prime labels,
\[
  \Gamma_c(X_{\rm site},\mathbb F_p^{2})
  \cong \mathbb F_p[\Lambda]^{2q}.
\]
This is the reason Haah's Pauli module is a free module over the translation
group algebra.  The basis elements \(g\in\Lambda\) label translated cells; they
are coefficient-basis labels, not prime ideals.

\subsection{The Sheaf-Theoretic Answer}

There is nevertheless a natural scheme in the background.  Once the
translation algebra is named,
\[
  R_{\rm tr}=\mathbb F_p[\Lambda],
\]
one can form
\[
  X_{\rm alg}=\Spec R_{\rm tr}.
\]
Haah's module complex can then be sheafified.  A generator module, Pauli module,
and stabilizer map
\[
  G\xrightarrow{\sigma}P
\]
give quasi-coherent sheaves
\[
  \widetilde G\xrightarrow{\widetilde\sigma}\widetilde P
\]
on \(X_{\rm alg}\).  With the translation-layer antipode, identity-coefficient
trace, and Laurent symplectic form fixed as in AQM-94, the excitation map gives
a complex
\[
  \widetilde G
  \xrightarrow{\widetilde\sigma}
  \widetilde P
  \xrightarrow{\widetilde\epsilon}
  \widetilde{G^\vee}.
\]
The algebraic diagnostic
\[
  \ker\epsilon/\operatorname{im}\sigma
\]
is therefore also a module, hence a quasi-coherent sheaf, when the quotient is
taken at the scheme-Haah layer of AQM-92.

This is the sheaf-theoretic reading that seems most faithful: Haah's
translation-invariant constraints are coherent module or sheaf data over the
spectral object \(X_{\rm alg}\).  But the physical tensor factors remain
indexed by \(X_{\rm site}\), and any identification between these two layers
requires an explicit bridge.

\subsection{Two Objects, Two Roles}

\begin{center}
\small\setlength{\tabcolsep}{3pt}
\begin{tabular}{p{0.24\linewidth}p{0.32\linewidth}p{0.34\linewidth}}
\textbf{Layer} & \textbf{Object} & \textbf{Role} \\
\hline
Site layer & \(\Lambda\times\{1,\ldots,q\}\), or
\(G\times\{1,\ldots,q\}\) periodically &
Indexes physical qubits or qudits and finite-support Pauli labels. \\
Translation algebra & \(R_{\rm tr}=\mathbb F_p[\Lambda]\) &
Packages translation shifts and finite-difference stabilizer matrices. \\
Scheme layer & \(X_{\rm alg}=\Spec R_{\rm tr}\) &
Supports sheaves \(\widetilde P,\widetilde G\), localization, fibres, support,
and algebraic defect modules. \\
Finite Weyl layer & finite quotient, residue, or Frobenius/self-dual data &
Turns selected labels into Weyl-Heisenberg operators and finite Hilbert
carriers. \\
\end{tabular}
\end{center}

The temptation is to collapse the first and third rows because Laurent
monomials look like coordinates.  AQM-91 separates this notation explicitly:
lowercase \(x,y\) are coordinate functions on an algebraic torus, while
uppercase \(X,Y\) are translation operators (AQM-91, lines 41--68).

\subsection{Example 1: Infinite Laurent Translation Algebra}

Let \(\Lambda=\mathbb Z^2\).  Then
\[
  R_{\rm tr}=\mathbb F_p[X^{\pm1},Y^{\pm1}].
\]
The physical site set is \(\mathbb Z^2\times\{1,\ldots,q\}\).  A finitely
supported Pauli label is a Laurent-polynomial vector in
\[
  P=R_{\rm tr}^{2q}.
\]
The sheaf-theoretic object is instead
\[
  \Spec\mathbb F_p[X^{\pm1},Y^{\pm1}],
\]
the split algebraic torus.  The exponent pair \((a,b)\) of \(X^aY^b\) is a
lattice translation label, not a point of this spectrum; AQM-94 records this
as Example 1 (AQM-94, lines 140--148).

\subsection{Example 2: Finite Periodic Quotient}

Let \(G=C_{N_x}\times C_{N_y}\).  The finite periodic translation algebra is
\[
  R_{G,p}
  =
  \mathbb F_p[X^{\pm1},Y^{\pm1}]
  /(X^{N_x}-1,Y^{N_y}-1).
\]
As a vector space it has one group-basis element \(X^iY^j\) for each periodic
site.  Multiplication by \(X\) and \(Y\) shifts those sites.  This is the
finite bridge used in AQM-91 (AQM-91, lines 70--91).

For the qubit toric square, AQM-91 writes
\[
  P_G=R_{G,2}^4
  \cong
  \mathbb F_2^{G\times\{x,y\}}{}_q
  \oplus
  \mathbb F_2^{G\times\{x,y\}}{}_p
\]
and recovers the Haah toric matrix from the Cayley-square boundary maps
(AQM-91, lines 93--138).  This is a successful finite site-to-module bridge.
It is still not the statement that the sites are the points of
\(\Spec R_{G,2}\).

\subsection{Example 3: Functions on a Finite Site Set}

For a finite set or finite group \(G\), the coordinate ring of the discrete
site scheme is
\[
  k^G=\operatorname{Map}(G,k).
\]
The scheme \(\Spec(k^G)\) has one reduced point per site when \(k\) is a
field, as in the product-field convention AQM-40.  This object is natural if
one wants functions on the site set.

The translation algebra \(k[G]\) is different.  It encodes convolution and
shifts.  A Haah finite-difference matrix uses \(k[G]\), not \(k^G\), unless a
Fourier or character decomposition has been explicitly named.  Thus a finite
model has two natural algebras:
\[
  k^G \quad\text{for functions on sites},\qquad
  k[G]\quad\text{for translations and shift operators}.
\]
They may be related in semisimple cases, but that relation is extra data.

\subsection{Example 4: Modular Thickening Warning}

Let \(G=C_3\) and work over \(\mathbb F_3\).  If \(g\) generates \(G\), then
\[
  \mathbb F_3[G]\cong \mathbb F_3[g]/(g^3-1)
  =\mathbb F_3[g]/((g-1)^3)
  \cong \mathbb F_3[u]/(u^3).
\]
As a translation module, the group basis is still \(\{1,g,g^2\}\), so a
periodic chain has three translated positions.  As a scheme, however,
\(\Spec\mathbb F_3[u]/(u^3)\) is one third-order thickened point.  AQM-94 uses
this as a warning: spectral points and translation-basis sites can differ
dramatically in modular characteristic (AQM-94, lines 174--184).

\subsection{Example 5: Split Character Basis Warning}

Over \(\mathbb F_3\),
\[
  \mathbb F_3[C_2]\cong \mathbb F_3[t]/(t^2-1)
  \cong \mathbb F_3\times\mathbb F_3.
\]
The spectrum has two reduced product-field points, while the translation basis
is \(\{1,t\}\).  Passing from the translation basis to the two idempotent
factors is a finite Fourier or character-basis step, not a tautological
identification of sites with \(\Spec\)-points.  AQM-94 records this split
product warning in its Example 5.

\subsection{Working Proposal}

The working proposal is to read Haah as a two-object construction.

\begin{enumerate}
\item The qubit or qudit tensor factors live over a site object
\(\Lambda\times\{1,\ldots,q\}\), or over a finite quotient
\(G\times\{1,\ldots,q\}\).
\item The translation-invariant stabilizer rules live as modules over
\(R_{\rm tr}=\mathbb F_p[\Lambda]\), or over a finite periodic quotient.
\item Those modules sheafify over \(\Spec R_{\rm tr}\), giving support,
localization, fibres, defect sheaves, and algebraic geometry.
\item A physical Weyl/Hilbert interpretation requires the translation-layer
finite-support convention, or a separate finite residue/Frobenius bridge.
\end{enumerate}

Thus Haah's definitions can be cleanly reinterpreted sheaf-theoretically, but
the sheaf geometry is not automatically the same object as the lattice of
qubits.  The bridge between the site object and the spectral object is part of
the data to construct, not a fact to assume.
